% Converted from README.md
\subsection{Valgrind}\label{valgrind}

\href{https://valgrind.org/}{Valgrind} je platforma za pravljenje alata za dinamičku analizu mašinskog koda, snimljenog ili kao objektni modul (nepovezan) ili kao izvršivi program (povezan). Postoje Valgrind alati koji mogu automatski da detektuju probleme sa memorijom i procesima.

Valgrind se može koristiti i kao alat za pravljenje novih alata. Valgrind distribucija, između ostalih, uključuje sledeće alate: detektor memorijskih grešaka (Memcheck), detektor grešaka u višenitnim programima (Helgrind i DRD), profiler keš memorije i grana (Cachegrind), profiler grafa poziva (Callgrind) i profiler korišćenja dinamičke memorije (Massif).

\subsubsection{Struktura i upotreba Valgrind alata}\label{struktura-i-upotreba-valgrind-alata}

Alat Valgrind se sastoji od alata za dinamičku analizu koda koji se kreira kao dodatak pisan u C programskom jeziku na jezgro Valgrinda. Jezgro Valgrinda omogućuje izvršavanje klijentskog programa, kao i snimanje izveštaja koji su nastali prilikom analize samog programa. Alati Valgrinda koriste metodu bojenja vrednosti. Oni svaki registar i memorijsku vrednost boje (zamenjuju) sa vrednošću koja govori nešto dodatno o originalnoj vrednosti. Proces rada svakog alata Valgrinda je u osnovi isti.

Valgrind deli originalni kod u sekvence koje se nazivaju osnovni blokovi. Osnovni blok je pravolinijska sekvenca mašinskog koda, na čiji se početak skače, a koja se završava skokom, pozivom funkcije ili povratkom u funkciju pozivaoca. Svaki kod programa koji se analizira ponovo se prevodi na zahtev, pojedinačno po osnovnim blokovima, neposredno pre izvršavanja samog bloka. Veličina osnovnog bloka je ograničena na maksimalno šezdeset mašinskih instrukcija.

Alat analizira dobijen kod i vrši translaciju - proces koji se sastoji od sledećih koraka:

\begin{enumerate}
\def\labelenumi{\arabic{enumi}.}
\tightlist
\item
  Disasembliranje (razgradnja) - prevođenje mašinskog koda u ekvivalentni interni skup instrukcija koje se nazivaju međukod instrukcije. U ovoj fazi međukod je predstavljen stablom. Ova faza je zavisna od arhitekture na kojoj se program izvršava.
\item
  Optimizacija 1 - prva faza optimizacije linearizuje prethodno izgrađeni međukod. Primenjuju se neke standardne optimizacije programskih prevodilaca kao što su uklanjanje redundantnog koda, eliminacija podizraza itd.
\item
  Instrumentacija - Blok međukoda se prosleduje alatu, koji može proizvoljno da ga transformiše. Prilikom instrumentacije alat u zadati blok dodaje dodatne međukod operacije, kojima proverava ispravnost rada programa. Treba napomenuti da ubačene instrukcije ne narušavaju konzistentno izvršavanje originalnog koda.
\item
  Optimizacija 2 - jednostavnija faza optimizacije od prve. Uključuje izračunavanje matematičkih izraza koji se mogu izvršiti pre faze izvršavanja i uklanjanje mrtvog koda.
\item
  Izgradnja stabla - linearizovani međukod se konvertuje natrag u stablo radi lakšeg izbora instrukcija.
\item
  Odabir instrukcija - stablo međukoda se konvertuje u listu instrukcija koje koriste virtuelne registre. Ova faza se takođe razlikuje u zavisnosti od arhitekture na kojoj se izvršava.
\item
  Alokacija registara - zamena virtuelnih registara stvarnim. Po potrebi se uvode prebacivanja u memoriju. Ne zavisi od platforme. Koristi se poziv funkcija koje pronalaze iz kojih se registara vrši čitanje i u koje se vrši upis.
\item
  Asembliranje - kodiranje izabranih instrukcija na odgovarajući način i smeštanje u blok memorije. Ova faza se takođe razlikuje u zavisnosti od arhitekture na kojoj se izvršava.
\end{enumerate}

Jezgro Valgrinda troši najviše vremena na sam proces pravljenja, pronalaženja i izvršavanja translacije (originalni kod se nikad ne izvršava). Treba napomenuti da sve ove korake osim instrumentacije izvršava jezgro Valgrinda, dok samu instrumentaciju izvršava određeni alat koji smo koristili za analizu izvornog koda.

Sve međukod instrukcije, originalne i dodate translacijom, prevode se u mašinske reči ciljne platforme i snimaju u prevedeni osnovni blok. Alat u originalni kod umeće operacije za instrumentaciju, zatim se takav kod prevodi.

Prilikom analize programa alatom Valgrind izvršavanje programa traje 20-100 puta duže nego inače. Analiza prevedenog programa Valgrindom, vrši sledećom naredbom:

\begin{minted}{bash}
valgrind --tool=alat [argumenti alata] ./a.out [argumenti za a.out]
\end{minted}

ili pokretanjem Valgrind memory analizer-a iz QtCreator-a za aktivan projekat.

Ukoliko se ne zada vrednost argumenta \texttt{-\/-tool} podrazumeva se \texttt{memcheck}.

Prve tri linije izlazne poruke štampaju se prilikom pokretanja bilo kog alata koji je u sklopu Valgrinda. U nastavku se prikazuju poruke o greškama koje je alat pronašao u programu. Zatim sledi izlaz samog programa, praćen sumiranim izveštajem o greškama.

Nekada informacija koja se dobije o grešci nije dovoljno detaljna da se u hiljadama linija koda nađe pravo mesto. Da bismo u okviru poruke o grešci imali i informaciju o liniji koda u kojoj je detektovana, potrebno je da program prevedemo sa simbolima za debagovanje (opcija \texttt{-g} za \texttt{gcc}). Da se ne bi dogodilo da se ne prijavljuje tačna linija u kojoj je detektovana greška, preporučuje se da se isključe optimizacije (opcija \texttt{-O0} za \texttt{gcc}). Nepotpuni stek okviri se manifestuju linijama koje sadrže nepoznate simbole (\texttt{???}), u kom slučaju se zarad njihovog poboljšanja, osim kompajlerskih opcija \texttt{-g} i \texttt{-O} mogu koristiti i \href{https://valgrind.org/docs/manual/faq.html\#faq.unhelpful}{neke Valgrind opcije}.
